\begin{solapartat}
\punts{0,75} per la representació correcta de les línies de camp. Si no s'indica el sentit de les línies, o les línies surten de la càrrega negativa o les línies convergeixen cap a la càrrega positiva es considerarà que la representació és totalment incorrecta. Si les línies no surten/entren radialment a prop de la càrrega puntual, es restarà 0,5 punts. Si les línies no estan distribuïdes homogèniament quan surten/entren a prop de la càrrega, es restarà 0,25 punts.

\punts{0,5} per la representació correcta de la projecció de les superfícies equipotencials. Si les línies equipotencials no són perpendiculars a les línies de camp es considerarà que la representació és totalment incorrecta.

No es penalitza no representar la línia que passa per l'origen. Cal tenir present que es demana una representació esquemàtica, no un gràfic acurat, per tant, no cal que la representació sigui molt acurada; el que es demana és que l'estudiant evidenciï que té clars els conceptes.

\figura[0.55]{sol_fig1.png}
\end{solapartat}

\begin{solapartat}
\punts{0,25} $E_p = E_n = k\dfrac{|Q|}{r^2}$, $E_p = E_n = 8,99\times 10^{9}\dfrac{1,5\times 10^{-6}}{0,05^2} = 5,39\times 10^{6}\ \mathrm{N/C}$

\punts{0,3} $\vec{E}_p = \vec{E}_n = 5,39\times 10^{6}\ \vec{\imath}\ \mathrm{N/C}$

\punts{0,3} $\vec{E} = \vec{E}_p + \vec{E}_n = 1,08\times 10^{7}\ \vec{\imath}\ \mathrm{N/C}$

\punts{0,4} $V_p = k\dfrac{|Q|}{r} = 2,70\times 10^{5}\ \mathrm{V}$, $V_n = -k\dfrac{|Q|}{r} = -2,70\times 10^{5}\ \mathrm{V}$

$V = V_p + V_n = 0\ \mathrm{V}$
\end{solapartat}
